% ============================================================================
\section{Capability Integration and Local Task Losses}
\label{sec:setting}
% ============================================================================

\subsection{Separately acquired capabilities as references}

Let $C_i(\theta)$ be the student's capability score on task $i$, with
larger values indicating better performance. Separately distilling each
teacher into the same student architecture gives reference scores $r_i$.
A capability-integration target is a shared student satisfying
\begin{equation}
    C_i(\theta)\geq r_i-\epsilon_i\quad\text{for every }i,
    \label{eq:capability_target}
\end{equation}
within a specified compute budget, where $\epsilon_i\geq0$ are tolerances
chosen before evaluating the joint run. The references measure separate
transfer; they do not establish that these targets are jointly attainable.
We report raw scores and compare learning curves at matched task exposure
and total compute. The local analysis below concerns distillation losses;
the experiments test how their changes relate to capability scores.

\subsection{Dense supervision on student rollouts}

There are $M$ tasks with prompt distributions $\mathcal D_i$ and assigned
frozen teachers $\pi_{T_i}$. The student and teachers share a vocabulary.
For a task-$i$ prompt $x$, the student generates a response $y$. At position
$u$, the teacher scores the student-generated prefix $c=(x,y_{<u})$.
Write $p_\theta(a\mid c)=\pi_\theta(a\mid c)$ and
$q_i(a\mid c)=\pi_{T_i}(a\mid c)$.

We use the teacher top-$k$ loss from MOPD
\citep{ma2026mopdmultiteacheronpolicydistillation}. Let
$S_i(c)=\operatorname{TopK}_k(q_i(\cdot\mid c))$. The position loss is
\begin{equation}
    \ell_i^{(k)}(\theta;c)=
    \sum_{a\in S_i(c)}
    \left[
        p_\theta(a\mid c)\log\frac{p_\theta(a\mid c)}{q_i(a\mid c)}
        -p_\theta(a\mid c)+q_i(a\mid c)
    \right].
    \label{eq:dense_topk_loss}
\end{equation}
The probabilities retain their full-vocabulary normalization. The
$-p_\theta+q_i$ correction makes the loss smallest when student and teacher
probabilities agree on the retained tokens. Both the teacher scores and the
generated prefixes are held fixed during the update. This objective uses
all retained token contributions at each observed position. For finite
$k$, its gradient is that of the specified surrogate; it need not equal
the full-vocabulary reverse-KL gradient.

Let $\mathcal R_i^t$ be the distribution of prompts and student responses
available for task $i$ at checkpoint $\theta_t$, conditional on training
history. We average positions within responses and responses within tasks:
\begin{equation}
    L_i^t(\theta)=
    \mathbb E_{(x,y)\sim\mathcal R_i^t}
    \left[\frac1{|y|}\sum_{u=1}^{|y|}
    \ell_i^{(k)}(\theta;x,y_{<u})\right].
    \label{eq:local_task_loss}
\end{equation}
The distribution $\mathcal R_i^t$ is held fixed when differentiating.
Positive weights $w_i$, fixed before training and summing to one, define
\begin{equation}
    F_t(\theta)=\sum_{i=1}^M w_iL_i^t(\theta).
    \label{eq:fixed_multitask_objective}
\end{equation}
Thus $\mu_i=\nabla L_i^t(\theta_t)$ and
$\mu_w=\sum_iw_i\mu_i=\nabla F_t(\theta_t)$ are the local mean
gradients. Our main comparison uses equal weights.

Updating the student also changes the distribution of generated states.
Adding and subtracting the loss on the previous distribution gives
\begin{equation}
\begin{split}
 L_i^{t+1}(\theta_{t+1})-L_i^t(\theta_t)
 ={}&L_i^t(\theta_{t+1})-L_i^t(\theta_t)\\
 &+L_i^{t+1}(\theta_{t+1})-L_i^t(\theta_{t+1}).
\end{split}
\label{eq:policy_drift}
\end{equation}
The first difference measures optimization on the previous states; the
second measures the change in the sampling distribution, including the
student's new rollouts. We therefore evaluate loss on fixed reference
prefixes separately from loss on fresh student responses. Neither loss
alone establishes the capability target in Equation~\ref{eq:capability_target}.

\subsection{Averaging task gradients before combining them}

A micro-batch contains $b$ prompts from one task. An optimizer step processes
$G$ micro-batches, with task $i$ receiving $m_i$ of them. The counts sum to
$G$ and satisfy $m_{\min}\leq m_i\leq m_{\max}$. We use
$m_{\min}\geq2$ to estimate each task's variability from the same step.

Let $g_{i,s}$ be the gradient of the response-mean loss in micro-batch $s$
of task $i$. Fixed-weight comparisons accumulate
\begin{equation}
    A=\sum_{i=1}^M w_i
       \underbrace{\frac1{m_i}\sum_{s=1}^{m_i}g_{i,s}}_{\text{task mean}}.
    \label{eq:batch_gradient_observation}
\end{equation}
Implementation only requires multiplying each micro-batch loss by $w_i/m_i$.
For example, an equally weighted task contributes one quarter of the
combined mean whether it receives two or six micro-batches; six
micro-batches give a more precise estimate of that mean.

Counts are chosen from history $\mathcal H_t$ before drawing the step's
data. If each fresh task-$i$ micro-batch has conditional mean $\mu_i$,
then $\mathbb E[A\mid\mathcal H_t,m]=\mu_w$. Changing counts preserves
this mean gradient. It can still change AdamW's expected update through
its moments and clipping. Response averaging removes the direct increase
in task weight caused by longer responses, one source of the length
imbalance studied by Open-MOPD
\citep{gao2026openmopddiagnosingfixingcapability}.
